% appendices/design_patterns.tex
\chapter{Design Patterns and Architectural Decisions}
\label{cha:design-patterns}

This appendix details the C++ design patterns and architectural decisions
underlying the SDR Control/Interface System and the RS-232 Interface,
expanding on the rationale only briefly mentioned in chapters
\autoref{cha:sdr-control} and \autoref{cha:rs232-interface}.


\section{Errors as Values}
\label{sec:design-patterns:errors}
Exceptions were not used in the SDR Control/Interface System or the RS-232
Interface, as they were deemed to introduce unnecessary overhead and implicit
error contracts. C-style error codes (e.g., returning an integer to indicate success or failure)
were also avoided, as they may lead to error-prone code due to forgetfulness on the part of the developer.
Instead, the design adopted a more modern approach using \texttt{std::expected<T,E>}, which allows functions to 
return either a successful result of type \texttt{T} or an error of type \texttt{E}, making error handling explicit and predictable, with the benefit of compile-time type checking.
This design choice promotes clearer and more maintainable code, as it forces the caller to handle potential errors explicitly, reducing the likelihood
of unhandled exceptions or overlooked error conditions.
Though it must be noted that this approach can lead to more verbose code at the call site, as it requires explicit checks for success or failure.


\section{Reuse of Standard Protocol Stacks}
\label{sec:design-patterns:reflection}
One of the most important architectural decisions made during the work was to
avoid inventing a project-specific application protocol for the serial link
once it became clear that the existing gRPC control surface already captured
the required operations well.

Instead, the serial channel was treated as a transport problem. By carrying
\gls{PPP} over \gls{RS232}, the system obtains a conventional point-to-point IP
link, after which the already existing gRPC interface can be reused unchanged.
This pushes responsibilities such as framing, packet transport, and network
bring-up into mature, standardized layers rather than custom project code.

The practical effect is a simpler architecture: fewer bespoke parsing rules,
less duplicated protocol logic, and better interoperability with standard
networking and diagnostic tooling.


\section{Transport-Agnostic Control Interface}
\label{sec:design-patterns:registry}
The final design treats the control API and the underlying transport as
separate concerns. The application-level contract is defined once through gRPC,
while the transport may be a conventional network connection or the
\gls{PPP}-backed serial path described in \autoref{cha:rs232-interface}.

This separation proved more valuable than maintaining different command formats
for different media. Once the lower layer exposes an IP channel, the upper
layer can remain unchanged, which reduces implementation effort and avoids
fragmenting the operational model seen by users and integrators.



\section{Iterations of the SDR Control System}
\label{sec:design-patterns:sdr-iterations}
 
The \gls{sdr-control-system} described in
\autoref{cha:sdr-control} is the third iteration of the system. The
two preceding attempts are summarised in
\autoref{tab:sdr-iterations}, as the reasoning behind their
abandonment directly informed several of the design decisions
discussed throughout this appendix.
 
\begin{table}[h]
\centering
\caption{Iterations of the SDR Control/Interface System.}
\label{tab:sdr-iterations}
\begin{tabular}{p{2.2cm} p{4.7cm} p{5.8cm}}
\toprule
\textbf{Iteration} & \textbf{Approach} & \textbf{Outcome / Reason for Abandonment} \\
\midrule
1 & 
    Reuse of an existing third-party program for \gls{SDR} control,
    adapted to the \gls{GNSS} use case rather than built from
    scratch. &
    Abandoned due to incompatibilities between the program's
    assumptions 
    (threading model, extensibility of its API) and the requirements of this
    internship, which made adaptation impractical. \\
\addlinespace
2 & A ground-up rewrite aiming for a fully general \gls{SDR} control
    library, supporting an open-ended range of modulation types
    and configurations. &
    Over-generalisation: most of the supported modulation types
    were never required by the actual use case, and the effort
    spent on generality (abstracting over modulation schemes that
    were never exercised) outweighed its benefit, while also
    making the codebase harder to reason about and test. \\
\addlinespace
3 (final) & A ground-up rewrite scoped specifically to the
    DTA-2115B/2116 cards and the \gls{GNSS} \gls{I/Q} transmission
    use case, described in \autoref{cha:sdr-control}. &
    Adopted. Narrowing the scope to the actual requirements
    made the system significantly simpler to implement, test, and
    reason about. \\
\bottomrule
\end{tabular}
\end{table}
 
This outcome is a useful counterpoint to the broader lesson drawn
in \autoref{cha:conclusion-retrospective}, where extensibility and
flexible interfaces are credited with absorbing the unplanned
integration effort described in \autoref{sec:scope-evolution}.
Taken together, the two findings suggest a more nuanced position
than either alone: designing for extensibility paid off where the
\emph{source} of change was external and unpredictable (integration
requirements imposed by GMV's pre-existing systems), but pursuing
generality \emph{pre-emptively}, ahead of any concrete requirement
for it, incurred a cost with no corresponding benefit in the second
iteration. The same tension resurfaces in
\autoref{sec:integration}, where the temptation to add another
special-case protocol for the serial path was ultimately rejected in favour of
reusing the already existing gRPC layer over a \gls{PPP}-backed link.
